\chapter{阶段结论与后续工作}
\section{已经验证的主要结果}
\subsection{公开动作策略提供了稳定的执行参照}
官方\pifive{}权重在四套LIBERO任务上完成117/120回合。图像置零或冻结后成功率大幅下降，说明该复现确实依赖视觉输入；执行旧轨迹或几何脚本不能替代这一结果。该部分属于公开方案复现及输入干预，不是自行训练出等同能力的通用动作模型。
\subsection{中文理解的收益需要与分层控制区分}
中文直接输入为22/30，冻结分层版本为28/30，仅翻译后交给平坦策略为29/30。分层没有超过仅翻译对照，后续固定切换和冻结视觉检查等实验也未显示稳定额外收益。正确的高层描述仍可能无法改变低层策略的目标偏好；这一接口问题比继续增大语言模型更直接。
\subsection{检查器改善离线判断，尚未改善整体任务成绩}
第二版检查器在近似完成样例上提高完成召回，但单次误报仍超过预设阈值。探索性闭环对照中，基座和适配版均为28/30；适配版正确确认次数反而较少。连续两票能挡住部分偶发错误，却不能保证相关误判被消除。里程碑记忆在单场景新初态的小批回合上改善停止决策，仍不足以证明跨技能通用完成判断。
\subsection{显式操作点在局部执行上产生了可测作用}
可供性低秩适配显著改善抽屉把手定位；同输入严格命中率在2B版本达到93.7\%。位置经深度反投影和规划后，在四个预接触回合上LoRA为4/4、基座为0/4、真值定位为4/4。这里验证的是到位，不是完整抓取。随后不重置状态交接给动作策略，完整任务仅2/4；原生策略对照揭示额外前置动作可能破坏策略熟悉的阶段分布。
\subsection{端侧部署完成了精度对照，速度仍不合格}
恢复后的Jetson具备GPU四位推理和适配器加载能力，并完成350帧、1014查询的PC与端侧对照。2B量化版本严格命中均为91.1\%，但逐条文字输出仅76.33\%完全一致。三次端侧视觉到PC预接触联合测试通过；15瓦模式下512像素输入的中位延迟5.062秒、第95百分位7.149秒，未达到单次3秒目标。缩小图像后的延迟变化有限，新的尺寸仍缺完整精度回归。
\section{对旧结果的修正}
显式目标适配分支后来发现低秩权重合并公式的实现错误。六组已保存输入中，旧部署与未合并参考的最大动作差约0.6055，修正后约0.01144。修正证明数值接口得到改善，但尚无修正后完整闭环复测。因此，旧低分不能直接解释成该类适配方法没有作用，修正后也不能预设成功率上升。

真实机械臂固定轨迹曾成功搬运物体，但重复时出现碰桌。它应作为数据和执行接口经验，而不是已经达到可靠自主抓取的依据。三条SmolVLA示教训练及受限试动没有形成稳定任务行为，50\%动作执行也未实际完成。仿真结果不用于替代这些真实硬件结论。
\section{从多个阶段得到的工程经验}
\begin{enumerate}
\item 先核对输入与动作语义，再解释模型能力。图像旋转、腕部偏置、归一化、动作块和低秩合并中的任一错误，都可能掩盖真正原因。
\item 为新增模块保留简单对照。翻译、固定切换、训练集中位点、原生策略和真值定位分别回答不同问题，不能只比较一条复杂链路的最终成功率。
\item 让评测与使用方式一致。检查器应看到与闭环一致的起始帧、当前帧和提示词；生成长度、量化范围和输入尺寸改变后也应重新检验精度。
\item 将定位、规划、到位、交接、任务成功和实时性分层验证。前一层通过能帮助定位后一层故障，却不意味着整条链路已通过。
\item 保留失败回合和版本条件。不同初态、重跑覆盖、测试集参与选择及理想深度等条件，会改变数字的解释。
\end{enumerate}
\section{后续工作按风险排序}
\subsection{恢复可信的动作适配对照}
首先在修正权重合并后的分支进行小规模数值一致性核对，再重做冻结初态的目标接口对照。只有确认部署动作与训练接口一致，才扩大训练数据或重新评价适配方案。
\subsection{扩展可供性和阶段兼容性}
保持原生策略参照，将操作点定位与阶段识别分开验证；优先测试相机变化、遮挡和新布局，明确何时需要规划预接触、何时应直接交给已有技能。跨场景、真实深度误差和多个技能仍是后续范围。
\subsection{端侧优化先定位瓶颈再改运行时}
现有剖析指向大量小算子和四位线性计算开销。后续优化应分别测输出长度、功率模式和运行时，但每次变更都保留同输入精度检查、峰值内存与交换空间记录。当前不将仿真同步等待服务视为实时达标。
\subsection{扩大独立验证而不是只增加演示}
当前若干关键比较仅有3--5回合，难以可靠估计尾部风险。下一轮应预先固定任务、初态、预算、检查点选择方法和失败归类，同时增加跨任务验证；用于调试的状态与正式测试分开。真实硬件验证需单独记录碰撞间隙、标定与停止条件，仿真通过后再开展。
